\subsection{FPTRANSA Common Model}\label{section:fptransa-common-model}
For an approximate transcendental instruction, $x'$ is the source after (:ref:FP.registers.STATE.FSTATUS.DAZ:), $y=f(x')$ is the
exact real reference value named by its accuracy contract, and $a$ is the approximate result before
(:ref:FP.registers.STATE.FSTATUS.FTZ:). For a format with precision $p$ and minimum normal exponent $e_{\min}$, define

\[
  \operatorname{ulp}_F(y) =
  \begin{cases}
    2^{\max(\lfloor\log_2 |y|\rfloor-p+1,\ e_{\min}-p+1)}, & y\ne0,\\
    2^{e_{\min}-p+1}, & y=0.
  \end{cases}
\]

The measured error is $|a-y|/\operatorname{ulp}_F(y)$. The bound applies only to finite reference values within
the selected format\textquotesingle{}s finite range; NaN, infinity, and overflow follow the instruction\textquotesingle{}s special-value rules.
Every present contract guarantees at most 4 ULP for both S and D, corresponding to at least 21 and 50 worst-case
relative precision bits respectively for nonzero normal results. CPUID may advertise a smaller certified bound.
Each instruction names its required \texttt{FPTRANSA\_ACCURACY} contract. A clear (:ref:FPTRANSA.cpuid.EXTENSIONS.FPTRANSA_ACCURACY.ACCURACY.PRESENT:) for that
contract raises (:ref:base.events.EXCEPTION.INVALID_OPERAND_RELATION:) before the instruction reads operands or floating-point state.
An approximate transcendental operation is available only when CPUID reports (:ref:FPTRANSA.cpuid.EXTENSIONS.DIRECTORY.FEATURES.FPTRANSA:) and marks that
instruction\textquotesingle{}s accuracy-contract ID present.

FPTRANSA applies DAZ before forming the exact reference input and measures its advertised ULP error before FTZ.
It ignores (:ref:FP.registers.STATE.FSTATUS.RM:) and formats its implementation-defined approximation with nearest-even rounding.
The same implementation, input, format, and relevant FSTATUS mode produce a deterministic approximation.
Different implementations may produce different bit patterns within the applicable accuracy contract.

A signaling NaN generates NV. NaN results are quieted and then processed by (:ref:FP.registers.STATE.FSTATUS.DN:).
(:ref:FP.registers.STATE.FSTATUS.FTZ:) replaces a subnormal result with the required signed zero and generates UF. Each instruction
generates NV and DZ according to its special-value rules, and generates OF and UF when required by its exact result;
NX is never generated, remains unchanged in FFLAGS, and cannot cause the (:ref:FP.events.EXCEPTION.FLOATING_POINT_EXCEPTION:)
architectural event. Overflow means that the exact result magnitude exceeds the selected
format\textquotesingle{}s maximum finite value. Underflow means that the exact nonzero result magnitude is below the selected format\textquotesingle{}s
minimum normal value. If any generated floating-point exception cause has its matching FSTATUS enable bit set, the
instruction raises the (:ref:FP.events.EXCEPTION.FLOATING_POINT_EXCEPTION:) architectural event before writing any destination. Otherwise
it commits its result or results and accrues the generated causes in FFLAGS as one architectural commit.

\clearpage
